\vspace*{40mm}
\begin{center}
{\Large\bfseries ON THE\par\medskip HANOVER TRIANGULATION.}
\end{center}
\clearpage\null\clearpage
\begin{center}
[I.]\par\medskip
{\Large\bfseries CORRESPONDENCE}\par\medskip
{\bfseries ON THE HANOVER TRIANGULATION.}\par\bigskip
[1.]\par\medskip
\textsc{Gauss} to \textsc{Schumacher}.\quad Göttingen, 5~July 1816.
\end{center}

\par\noindent\hspace*{\parindent}First of all, my heartfelt congratulations on the splendid undertaking on such a grand scale that you announce to me in your last letter. This arc measurement in the royal Danish states will in itself tell us a great deal about the shape of the Earth. I have no doubt, however, that it will be possible in the future to extend your measurements southward through the Kingdom of Hanover. At the moment I cannot yet raise such a proposal in Hanover, since astronomy itself still needs so much support. But I am convinced that our government, which is also glad to support the sciences, will before long follow the glorious example of your excellent king. We would then already have a respectable meridian arc of \(6\frac{1}{2}\) degrees, and it would also be easy to link these operations to the Bavarian triangles. The latter were certainly measured with the greatest care, and it is regrettable that they are being withheld from publication.

\par\noindent\hspace*{\parindent}I have devised my own method of treating the measured triangles in the calculation, but it would take far too long to explain in a letter. If I have yet to publish it by the time you have measured your triangles, I shall discuss it with you in detail. Indeed, I offer to undertake the calculation of the principal triangles myself.
\par\medskip\noindent IX.\hfill44
\clearpage
\noindent346\hfill{\small\textbf{CORRESPONDENCE WITH SCHUMACHER.}}\par\medskip

\par\noindent\hspace*{\parindent}I have only one slight doubt about the second part of your undertaking, the measurement of a degree of longitude. My impression was that the lands of the Danish monarchy could be described as rather \textls[120]{flat}, or at least that they have no high mountains. If that assumption is sound and it consequently compels you to use an intermediate point, or even \textls[120]{several}, to determine the difference in astronomical longitude, some slight uncertainty will always attach to that determination, however skilled your assistants and however good your equipment.

\par\noindent\hspace*{\parindent}You have a great advantage in the fact that Denmark has already been surveyed trigonometrically. The measured angles themselves, of course, fall far short of what is needed for an arc measurement. It is those earlier operations that will make it much easier for you to choose your stations. Finding the stations would be the most disagreeable part of a similar undertaking for me, precisely because it wastes so much time. I have gone to great lengths (with similar future operations in mind) to obtain the angles measured by \textsc{Epailly} in Hanover, but without success.

\par\noindent\hspace*{\parindent}.\ .\ .\ .\ .\ .\quad An interesting problem occurred to me, namely
\begin{quote}
\noindent»\textbf{in general}, to project (map) a given surface onto another (given) surface in such a way that the image is similar to the original in the smallest parts.«
\end{quote}
\par\noindent\hspace*{\parindent}A special case is when the first surface is a sphere and the second a plane. Here the stereographic and Mercator projections are particular solutions. What one seeks, however, is the \textls[120]{general} solution for every kind of surface, encompassing all particular solutions.

\par\noindent\hspace*{\parindent}I understand that \textsc{Monge} and \textsc{Poinsot} have already worked on this in the Journal philomathique (as \textsc{Burckhardt} wrote to \textsc{Lindenau}). But since I do not know precisely \textls[120]{where}, I have not yet been able to look it up and therefore do not know whether those gentlemen's solutions correspond exactly to my idea and exhaust the subject. If they do not, this would seem to me a suitable prize question for a society to set someday.\quad.\ .\ .\ .\ .\ .
